\RequirePackage{silence}
\WarningsOff % Turn off all warnings from all packages
\ErrorsOff   % (Caution) Use only if you are sure the code has no severe logic errors
\documentclass[a4paper,12pt, twoside]{memoir}
\usepackage[utf8x]{inputenc}
\usepackage[english]{babel}    
\usepackage{url}
\usepackage{microtype}
\usepackage{color}
\usepackage{amsmath}
\usepackage{amsthm}         % Provides proof environment
\usepackage{amssymb}
\usepackage{mathtools}     % Includes amsmath but optimized and with more features
\usepackage{hyperref}
\usepackage{enumitem}
\usepackage{autobreak}
\usepackage{array}          % Support for notation table formatting
\usepackage{booktabs}       % Professional table rules

\hypersetup{
	colorlinks=true,
	linkcolor=blue,
	filecolor=magenta,
	urlcolor=cyan
}

% Theorem Environments
\newtheorem{definition}{Definition}[chapter]
\newtheorem{md}{Proposition}[chapter]
\newtheorem{theorem}{Theorem}[chapter]
\newtheorem{proposition}{Proposition}[chapter]
\newtheorem{corollary}{Corollary}[chapter]
\newtheorem{remark}{Remark}[chapter]
\newtheorem{example}{Example}[chapter]

% Declare Math Operators
\DeclareMathOperator{\Tr}{Tr}
\DeclareMathOperator{\Sym}{Sym}
\DeclareMathOperator{\setu}{s}
\DeclareMathOperator{\dimU}{\dim}
\DeclareMathOperator{\w}{w}
\DeclareMathOperator{\card}{card}
\DeclareMathOperator{\ucard}{ucard}
\DeclareMathOperator{\UCARD}{ucard}
\DeclareMathOperator{\pr}{pr}
\DeclareMathOperator{\pgr}{pgr}
\DeclareMathOperator{\scale}{scale}
\DeclareMathOperator{\rad}{rad}
\DeclareMathOperator{\Conf}{Conf}

% Custom commands
\newcommand{\set}[1]{\left\{#1\right\}} % set
\newcommand{\tbd}[1]{\mathcal{#1}} % uncertain set (styled uppercase)
\newcommand{\s}[1]{\set{#1}_\mathrm{u}} % uncertain number
\newcommand{\kbd}[1]{{#1}_\mathrm{u}} % uncertain interval
\newcommand{\bd}[1]{\left(#1\right)_\mathrm{u}} % uncertain expression
\newcommand{\tsbd}[1]{\mathbf{U}(\mathbb{#1})}
\newcommand{\tsbdcapn}[2]{\mathbf{U^{#2}}(\mathbb{#1})}
\newcommand{\tsbdts}[1]{\mathbf{U}_{\w}(\mathbb{#1})} % set of uncertain numbers
\newcommand{\ddclmd}{\mathcal{V}} % truth measure of a proposition
\newcommand{\ts}[1]{\w_{#1}} % Weight
\newcommand{\gtcl}[2]{\left(#1\right)_{#2}} % truth measure
\newcommand{\gr}[1]{\text{gr}(#1)} % graph of a function
\newcommand{\hpw}[2]{#1_1\left(#2\right)} % point-wise function
\newcommand{\hbd}[2]{#1_m\left(#2\right)} % uncertain function
\newcommand{\eq}[1]{\left[#1\right]_{eq}} % eqs operator
\newcommand{\kgdd}[1]{\left(#1\right)_1} % pointwise operator
\newcommand{\kgddmr}[1]{\left(#1\right)_{1'}} % real operator
\newcommand{\kgdt}[1]{\left(#1\right)_m} % pointwise operator
\newcommand{\kgdtmr}[1]{\left(#1\right)_{m'}} % pointwise operator

\title{Scenario-Based Uncertain Set Theory}
\author{
	Tran Manh Cuong\thanks{Email: \href{mailto:tmcuong2408@gmail.com}{tmcuong2408@gmail.com} | Phone: 0353237140} \\[0.5em]
	\small Alumnus of Faculty of Mathematics and Informatics (K11)
	\\ Thai Nguyen University of Sciences \\
	\small Address: DJ7 Street, Thoi Hoa, Ho Chi Minh City, Vietnam \\
	\small \textit{Research Field: Mathematics of Uncertainty}
}
\date{Last updated: \today}

\begin{document}
	
	\maketitle
	
	\begin{abstract}
		This paper introduces Scenario-Based Uncertain Set Theory, a foundational mathematical framework built from first principles to model structural uncertainty without relying on a priori probability distributions, fuzzy membership grades, or static interval bounds \cite{tran2026extended}. Unlike classical Interval Arithmetic—which frequently suffers from dependency loss and numerical overestimation—our approach represents Uncertain Elements directly as non-empty possibility domains $X \subseteq U$ and constructs Uncertain Sets as styled-uppercase indexed families $\tbd{S} = \{X_i : i \in I\}$.
		
		Through this scenario-based formulation, every uncertain set $\tbd{S}$ canonically induces a Configuration Space $\mathrm{Conf}(\tbd{S})$ of realized classical sets. Consequently, set cardinality naturally generalizes to an uncertain number $\UCARD(\tbd{S}) \subseteq \mathbb{N}$, seamlessly capturing state-collapse phenomena when distinct elements realize identical values. By equipping the global Cartesian scenario space $\Omega_{\tbd{S}}$ with a normalized product measure, we establish a rigorous truth valuation $\mathcal{V}(u \in_u \tbd{S}) \in [0, 1]$ for element membership and prove the Conservation of Total Membership Valuation theorem, showing that $\sum_{u \in U} \mathcal{V}(u \in_u \tbd{S}) = \mathbb{E}[\UCARD(\tbd{S})]$. Finally, we extend the framework to Uncertain Mappings $f : \tbd{A} \to \tbd{B}$, demonstrating how fiber-based uncertain lifting resolves non-injective obstructions in classical functional equations.
		
		\vspace{1em}
		\noindent \textbf{Keywords:} Uncertain Set Theory, Uncertain Elements, Configuration Space, Uncertain Cardinality, Scenario Space, Membership Valuation, Uncertain Mapping, First Principles.
	\end{abstract}
	
	\chapter{Uncertain Sets}
	
	\begin{definition}[Uncertain Element]
		Let $U$ be a non-empty universal set. An uncertain element $X$ on $U$ is an object whose exact value in $U$ is undetermined, identified directly with its non-empty set of possible values (possibility domain or state space):
		\[
		X \subseteq U, \quad X \neq \emptyset
		\]
		Each realized state of $X$ is a classical element $v \in X$.
	\end{definition}
	
	\begin{definition}[Uncertain Set]
		Let $U$ be a non-empty universal set and $I$ be an index set. An uncertain set $\tbd{S}$ on $U$ is an indexed family of uncertain elements on $U$:
		\[
		\tbd{S} = \{X_i : i \in I\}
		\]
		where each uncertain element $X_i \subseteq U$ ($i \in I$) is a non-empty possibility domain.
	\end{definition}
	
	\begin{definition}[Configuration Space]
		Let $\tbd{S} = \{X_i : i \in I\}$ be an uncertain set on $U$. A realized configuration (or empirical state) of $\tbd{S}$ is a classical set $S^* \subseteq U$ obtained when all constituent uncertain elements $X_i$ simultaneously realize specific values $v_i \in X_i$:
		\[
		S^* = \{v_i : i \in I\}
		\]
		The collection of all classical sets realized by $\tbd{S}$ is called the configuration space of $\tbd{S}$, denoted by $\mathrm{Conf}(\tbd{S})$.
	\end{definition}
	
	\begin{definition}[Classical Cardinality Operator]
		Let $A$ be a classical set. The classical cardinality of $A$, denoted by $\mathrm{card}(A)$ or $|A|$, measures the number of distinct elements in $A$. For any configuration $S^* \in \mathrm{Conf}(\tbd{S})$, $\mathrm{card}(S^*)$ maps $S^*$ to its cardinal value in $\mathbb{N} \cup \{\infty\}$.
	\end{definition}
	
	\begin{definition}[Uncertain Cardinality]
		Let $\tbd{S} = \{X_i : i \in I\}$ be an uncertain set on $U$. The uncertain cardinality (or size) of $\tbd{S}$, denoted by $\UCARD(\tbd{S})$, is an uncertain number on $\mathbb{N}$ consisting of the classical cardinalities of all realized states in its configuration space $\mathrm{Conf}(\tbd{S})$:
		\[
		\UCARD(\tbd{S}) = \{ \mathrm{card}(S^*) : S^* \in \mathrm{Conf}(\tbd{S}) \} \subseteq \mathbb{N}
		\]
		where $\mathrm{card}(S^*) = |S^*|$ represents the classical cardinality of the configuration $S^*$.
	\end{definition}
	
	\begin{remark}
		If all uncertain elements $X_i \in \tbd{S}$ take mutually distinct values across all realized configurations, then $|S^*| = |I|$ for every $S^* \in \mathrm{Conf}(\tbd{S})$, in which case the uncertain cardinality degenerates into a classical singleton constant $\UCARD(\tbd{S}) = \{|I|\}$. Conversely, cardinality degeneration (state collapse) occurs at configurations $S^*$ where duplicate values appear ($v_i = v_j$ for some $i \neq j$).
	\end{remark}
	
	\begin{example}[Numerical Illustration of Configuration Space and Uncertain Cardinality]
		Consider the universal set $U = \mathbb{R}$. Let $X_1$ and $X_2$ be two uncertain elements given by:
		\[
		X_1 = \{1, 2\}, \quad X_2 = \{1, 3\}
		\]
		Let $\tbd{S} = \{X_1, X_2\}$ be an uncertain set.
		
		Assuming $X_1$ and $X_2$ vary independently, the joint scenario space yields $2 \times 2 = 4$ realized classical configurations:
		\begin{itemize}
			\item $v_1 = 1, v_2 = 1 \implies S_1^* = \{1, 1\} = \{1\}$, with $|S_1^*| = \mathrm{card}(S_1^*) = 1$.
			\item $v_1 = 1, v_2 = 3 \implies S_2^* = \{1, 3\}$, with $|S_2^*| = \mathrm{card}(S_2^*) = 2$.
			\item $v_1 = 2, v_2 = 1 \implies S_3^* = \{2, 1\}$, with $|S_3^*| = \mathrm{card}(S_3^*) = 2$.
			\item $v_1 = 2, v_2 = 3 \implies S_4^* = \{2, 3\}$, with $|S_4^*| = \mathrm{card}(S_4^*) = 2$.
		\end{itemize}
		
		Thus, the configuration space of $\tbd{S}$ is explicitly:
		\[
		\mathrm{Conf}(\tbd{S}) = \big\{ \{1\}, \{1, 3\}, \{1, 2\}, \{2, 3\} \big\}
		\]
		The uncertain cardinality $\UCARD(\tbd{S})$ is the uncertain number on $\mathbb{N}$ given by:
		\[
		\UCARD(\tbd{S}) = \{ |S^*| : S^* \in \mathrm{Conf}(\tbd{S}) \} = \{1, 2\} \subseteq \mathbb{N}
		\]
	\end{example}
	
	\begin{definition}[Element Membership Valuation via Scenario Measure]
		Let $\tbd{S} = \{X_1, X_2, \dots, X_n\}$ be an uncertain set on a non-empty universe $U$, where each uncertain element $X_i \subseteq U$ varies independently.
		
		The global scenario space of $\tbd{S}$ is the Cartesian product:
		\[
		\Omega_{\tbd{S}} = X_1 \times X_2 \times \dots \times X_n
		\]
		equipped with a product measure $\lambda$. Each state tuple $\mathbf{v} = (v_1, \dots, v_n) \in \Omega_{\tbd{S}}$ uniquely induces a realized configuration $S^*(\mathbf{v}) = \{v_1, \dots, v_n\} \in \mathrm{Conf}(\tbd{S})$.
		
		For any classical element $u \in U$, let $E_u \subseteq \Omega_{\tbd{S}}$ be the measurable set of scenarios whose realized configuration contains $u$:
		\[
		E_u = \{\mathbf{v} \in \Omega_{\tbd{S}} \;:\; u \in S^*(\mathbf{v})\}
		\]
		The truth value of the membership assertion $u \in_u \tbd{S}$ is defined by the normalized measure:
		\[
		\mathcal{V}(u \in_u \tbd{S}) := \frac{\lambda(E_u)}{\lambda(\Omega_{\tbd{S}})} \in [0, 1]
		\]
	\end{definition}
	
	\begin{example}[Numerical Illustration of Membership Valuation]
		Consider $U = \mathbb{R}$ and independent uncertain elements $X_1, X_2$:
		\[
		X_1 = \{1, 3, 5, 7\}, \quad X_2 = \{2, 3, 4, 5\}
		\]
		Let $\tbd{S} = \{X_1, X_2\}$. The joint scenario space is $\Omega_{\tbd{S}} = X_1 \times X_2$. Equipping $\Omega_{\tbd{S}}$ with the counting measure $\lambda$, the total measure is $\lambda(\Omega_{\tbd{S}}) = 4 \times 4 = 16$.
		
		The 16 state tuples $\mathbf{v} = (v_1, v_2) \in \Omega_{\tbd{S}}$ and their realized configurations $S^*(\mathbf{v})$ are:
		\begin{align*}
			\Omega_{\tbd{S}} = \{ & (1,2), (1,3), (1,4), (1,5), (3,2), (3,3), \\
			& (3,4), (3,5), (5,2), (5,3), (5,4), (5,5), \\
			& (7,2), (7,3), (7,4), (7,5) \}
		\end{align*}
		
		To evaluate the truth measure of $3 \in_u \tbd{S}$, we identify $E_3 \subseteq \Omega_{\tbd{S}}$ where $3 \in S^*(\mathbf{v})$:
		\[
		E_3 = \{(1,3), (3,2), (3,3), (3,4), (3,5), (5,3), (7,3)\}
		\]
		Since $\lambda(E_3) = |E_3| = 7$, the truth valuation is:
		\[
		\mathcal{V}(3 \in_u \tbd{S}) = \frac{\lambda(E_3)}{\lambda(\Omega_{\tbd{S}})} = \frac{7}{16} = 0.4375
		\]
	\end{example}
	
	\begin{theorem}[Conservation of Total Membership Valuation]
		Let $U$ be a non-empty universal set and let $\tbd{S} = \{X_1, X_2, \dots, X_n\}$ be an uncertain set consisting of $n$ finite uncertain elements $X_i \subseteq U$.
		
		Then, the sum of the membership truth valuations of all elements $u \in U$ equals the expected value of the uncertain cardinality $\UCARD(\tbd{S})$ over the finite scenario space $\Omega_{\tbd{S}}$:
		\[
		\sum_{u \in U} \mathcal{V}(u \in_u \tbd{S}) = \mathbb{E}\left[\UCARD(\tbd{S})\right]
		\]
		where the expected cardinality $\mathbb{E}\left[\UCARD(\tbd{S})\right]$ is defined as the arithmetic mean of realized cardinalities across all scenarios:
		\[
		\mathbb{E}\left[\UCARD(\tbd{S})\right] = \frac{1}{|\Omega_{\tbd{S}}|} \sum_{\mathbf{v} \in \Omega_{\tbd{S}}} |S^*(\mathbf{v})|
		\]
		and $|S^*(\mathbf{v})|$ represents the classical cardinality of the realized configuration $S^*(\mathbf{v}) = \{v_1, \dots, v_n\}$ at scenario $\mathbf{v} = (v_1, \dots, v_n) \in \Omega_{\tbd{S}}$.
	\end{theorem}
	
	\begin{proof}
		By Definition 1.6, for each element $u \in U$, the truth valuation of $u \in_u \tbd{S}$ under the counting measure is:
		\[
		\mathcal{V}(u \in_u \tbd{S}) = \frac{|E_u|}{|\Omega_{\tbd{S}}|}
		\]
		where $E_u = \{\mathbf{v} \in \Omega_{\tbd{S}} \;:\; u \in S^*(\mathbf{v})\}$ is the subset of scenarios whose realized configuration contains $u$.
		
		Using the indicator function $\mathbb{I}_{\{u \in S^*(\mathbf{v})\}}(\mathbf{v})$, the cardinality $|E_u|$ is expressed as:
		\[
		|E_u| = \sum_{\mathbf{v} \in \Omega_{\tbd{S}}} \mathbb{I}_{\{u \in S^*(\mathbf{v})\}}(\mathbf{v})
		\]
		
		Summing the truth valuations over all $u \in U$:
		\[
		\sum_{u \in U} \mathcal{V}(u \in_u \tbd{S}) = \frac{1}{|\Omega_{\tbd{S}}|} \sum_{u \in U} \sum_{\mathbf{v} \in \Omega_{\tbd{S}}} \mathbb{I}_{\{u \in S^*(\mathbf{v})\}}(\mathbf{v})
		\]
		
		Since $S^*(\mathbf{v})$ contains at most $n$ elements for any scenario $\mathbf{v}$, the set of elements $u \in U$ with non-zero indicator is finite. Interchanging the order of double summation yields:
		\[
		\sum_{u \in U} \sum_{\mathbf{v} \in \Omega_{\tbd{S}}} \mathbb{I}_{\{u \in S^*(\mathbf{v})\}}(\mathbf{v}) = \sum_{\mathbf{v} \in \Omega_{\tbd{S}}} \left( \sum_{u \in U} \mathbb{I}_{\{u \in S^*(\mathbf{v})\}}(\mathbf{v}) \right)
		\]
		
		For a fixed scenario tuple $\mathbf{v} \in \Omega_{\tbd{S}}$, summing the indicators over all $u \in U$ counts exactly the number of distinct elements in the realized configuration $S^*(\mathbf{v})$:
		\[
		\sum_{u \in U} \mathbb{I}_{\{u \in S^*(\mathbf{v})\}}(\mathbf{v}) = |S^*(\mathbf{v})|
		\]
		
		Substituting this back into the summation completes the proof:
		\[
		\sum_{u \in U} \mathcal{V}(u \in_u \tbd{S}) = \frac{1}{|\Omega_{\tbd{S}}|} \sum_{\mathbf{v} \in \Omega_{\tbd{S}}} |S^*(\mathbf{v})| = \mathbb{E}\left[\UCARD(\tbd{S})\right]
		\]
	\end{proof}
	\chapter{Uncertain Mapping}
	\begin{definition}[Uncertain Mapping]
		Let $\tbd{A}$ and $\tbd{B}$ be two non-empty uncertain sets on universal sets $U$ and $V$, respectively. An uncertain mapping (or uncertain function) $f$ from $\tbd{A}$ to $\tbd{B}$, denoted by $f : \tbd{A} \to \tbd{B}$, is a rule that assigns to each uncertain element $X \in \tbd{A}$ a unique uncertain element $Y = f(X) \in \tbd{B}$.
		Here, $\tbd{A}$ is called the {domain} of $f$, $\tbd{B}$ is called the {codomain} of $f$, and $f(X)$ is called the {image} of $X$ under $f$.
	\end{definition}
	\begin{example}[Independent Coin Tossing Mapping]
		Consider tossing a fair coin twice, where $1$ represents heads and $0$ represents tails. Let $X = \{0, 1\}$ be the uncertain element representing a single toss, and let $\tbd{A} = \{X\}$ and $\tbd{B} = \{Y\}$ be uncertain sets on $\mathbb{R}$.
		
		Define the uncertain mapping $f_1 : \tbd{A} \to \tbd{B}$ that maps $X \in \tbd{A}$ to $Y = f_1(X) \in \tbd{B}$, where $Y$ is defined via the Minkowski sum:
		\[
		Y = \{v_1 + v_2 : v_1, v_2 \in X\} = \{0, 1, 2\}
		\]
		Thus, $f_1$ maps the single-toss uncertain element $X$ to the two-toss total uncertain element $Y = \{0, 1, 2\}$.
	\end{example}
	
	\begin{example}[Scaled Single Coin Tossing Mapping]
		Consider tossing a fair coin once, where $1$ represents heads and $0$ represents tails, and doubling the resulting score. Let $X = \{0, 1\}$ be the uncertain element representing the single toss, and let $\tbd{A} = \{X\}$ and $\tbd{B} = \{Y\}$ be uncertain sets on $\mathbb{R}$.
		
		Define the uncertain mapping $f_2 : \tbd{A} \to \tbd{B}$ that maps $X \in \tbd{A}$ to $Y = f_2(X) \in \tbd{B}$, where $Y$ is defined via:
		\[
		Y = \{2v : v \in X\} = \{0, 2\}
		\]
		Thus, $f_2$ maps the single-toss uncertain element $X$ to the doubled-score uncertain element $Y = \{0, 2\}$.
	\end{example}
	\begin{example}[Functional Equation Induced Uncertain Mapping]
		Consider a formal classical functional relation $f : [-1, 1] \to [-1, 1]$ satisfying the trigonometric identity:
		\[
		f(\sin x) = \cos x \quad \text{for all } x \in \mathbb{R}
		\]
		In classical analysis, no single-valued function $f : [-1, 1] \to [-1, 1]$ satisfies this relation because $\sin x$ is non-injective on $\mathbb{R}$ (for instance, $x = \frac{\pi}{6}$ yields $f\left(\frac{1}{2}\right) = \frac{\sqrt{3}}{2}$, whereas $x = \frac{5\pi}{6}$ yields $f\left(\frac{1}{2}\right) = -\frac{\sqrt{3}}{2}$).
		
		To resolve this obstruction within Uncertain Set Theory, for each parameter $t \in [-1, 1]$, let $X_t$ be the uncertain element on $U = \mathbb{R}$ representing the root-fiber of $\sin v = t$:
		\[
		X_t = \{ v \in \mathbb{R} : \sin v = t \} = \{ \arcsin t + 2k\pi, \; \pi - \arcsin t + 2k\pi : k \in \mathbb{Z} \}
		\]
		Since the correspondence $t \mapsto X_t$ forms a canonical bijection between $[-1, 1]$ and the family of root-fiber uncertain elements $\tbd{X} = \{ X_t : t \in [-1, 1] \}$, we establish the contextual identification $t \equiv X_t$ via the embedding operator $\iota : [-1, 1] \to \tbd{X}$ defined by $\iota(t) = X_t$.
		
		Let $\tbd{Y} = \{ Y_t : t \in [-1, 1] \}$ be an uncertain set on $V = [-1, 1]$, where each $Y_t$ is the uncertain image element corresponding to $X_t$:
		\[
		Y_t = \{ \cos v : v \in X_t \} = \left\{ -\sqrt{1 - t^2}, \; \sqrt{1 - t^2} \right\}
		\]
		Then, the rule $F : \tbd{X} \to \tbd{Y}$ assigning each $X_t \in \tbd{X}$ to $Y_t = F(X_t) \in \tbd{Y}$ defines a valid single-valued uncertain mapping satisfying the lifted identity:
		\[
		F\big(\iota(\sin X_t)\big) = \cos X_t \quad \text{for all } t \in [-1, 1]
		\]
		In particular, for $t = \frac{1}{2}$, $F$ maps the countably infinite root-set uncertain element $X_{1/2}$ to the dual-valued uncertain element $Y_{1/2} = \left\{ -\frac{\sqrt{3}}{2}, \; \frac{\sqrt{3}}{2} \right\}$, whereas at the boundary points $t = \pm 1$, $Y_{\pm 1}$ degenerates to the classical singleton constant $\{0\}$.
	\end{example}
	\begin{thebibliography}{9}
		\bibitem{tran2026extended}
		Tran, M. C. (2026). \textit{Extended Logic and Mathematics of Uncertainty}. Zenodo. DOI: \href{https://doi.org/10.5281/zenodo.22031881}{10.5281/zenodo.22031881}.
		\bibitem{mu}
		tmcuong2408. (2026). tmcuong2408/mu: \textit{Mathematics of Uncertainty Library (Version v0.2.0)} [Computer software]. Zenodo.
		DOI: \href{https://doi.org/10.5281/zenodo.22051784}{10.5281/zenodo.22051784}
	\end{thebibliography}
	
\end{document}